% Part: methods
% Chapter: induction
% Section: relations

\documentclass[../../../include/open-logic-section]{subfiles}

\begin{document}

\olfileid{mth}{ind}{rel}

\olsection{संबंध आणि फलने}

नैसर्गिक संख्या किंवा नीटस पदे यांसारख्या वस्तूंचा संच विगमनाधारित रीतीने परिभाषित केल्यावर, त्या वस्तूंवरील \emph{संबंध}ही विगमनाने परिभाषित करता येतात. उदाहरणार्थ, एका नीटस पद~$t_2$ मध्ये दुसरे नीटस पद~$t_1$ भाग म्हणून येत असेल, तर पहिले पद दुसऱ्याचे उपपद आहे, अशी कल्पना करा. त्यासाठी $t_1 \sqsubseteq t_2$ हे चिन्ह वापरू. अर्थात, प्रत्येक नीटस पद स्वतःचे उपपद असते: $t \sqsubseteq t$. या संबंधाची विगमनाधारित व्याख्या अशी देता येते:

\begin{defn}
नीटस पद~$t_1$ हे~$t_2$ चे उपपद आहे हा संबंध, $t_1 \sqsubseteq t_2$, पुढीलप्रमाणे~$t_2$ वर विगमनाने परिभाषित केला आहे:
\begin{enumerate}
\item $t_2$ हे अक्षर असेल, तर $t_1 \sqsubseteq t_2$ हे $t_1 = t_2$ असेल तेव्हाच आणि फक्त तेव्हाच सत्य असते.
\item $t_2$ हे $[s_1 \circ s_2]$ असेल, तर $t_1 \sqsubseteq t_2$ हे $t_1 = t_2$, $t_1 \sqsubseteq s_1$ किंवा $t_1 \sqsubseteq s_2$ यांपैकी एक असेल तेव्हाच आणि फक्त तेव्हाच सत्य असते.
\end{enumerate}
\end{defn}

उदाहरणार्थ, या व्याख्येनुसार $\mathrm{a} \sqsubseteq [\mathrm{b} \circ \mathrm{a}]$ मिळते. कारण (2) नुसार $\mathrm{a} \sqsubseteq [\mathrm{b} \circ \mathrm{a}]$ हे $\mathrm{a} = [\mathrm{b} \circ \mathrm{a}]$, किंवा $\mathrm{a} \sqsubseteq b$, किंवा $\mathrm{a} \sqsubseteq \mathrm{a}$ यांपैकी एक असेल तेव्हाच आणि फक्त तेव्हाच सत्य असते. पहिले दोन पर्याय असत्य आहेत: $\mathrm{a}$ हे स्पष्टपणे $[\mathrm{b} \circ \mathrm{a}]$ शी एकरूप नाही; आणि (1) नुसार $\mathrm{a} \sqsubseteq \mathrm{b}$ हे $\mathrm{a} = \mathrm{b}$ असेल तेव्हाच आणि फक्त तेव्हाच सत्य असते; तेही असत्य आहे. पण पुन्हा (1) नुसार $\mathrm{a} \sqsubseteq \mathrm{a}$ हे $\mathrm{a} = \mathrm{a}$ असेल तेव्हाच आणि फक्त तेव्हाच सत्य असते; ते सत्य आहे.

ही व्याख्या यशस्वी होण्यासाठी एक अजून न सिद्ध केलेले तथ्य आवश्यक आहे: प्रत्येक नीटस पद~$t$ हे स्वतंत्र अक्षर असते, किंवा अशी \emph{एकमेव रीतीने निश्चित} होणारी नीटस पदे $s_1$ आणि $s_2$ असतात की $t = [s_1 \circ s_2]$. येथे “एकमेव रीतीने निश्चित” याचा अर्थ असा: $t = [s_1 \circ s_2]$ असेल, तर $s_1 \neq r_1$ किंवा $s_2 \neq r_2$ असताना ते \emph{त्याचबरोबर} $= [r_1 \circ r_2]$ असू शकत नाही. असे घडले, तर खंड~(2) चाच स्वतःशी संघर्ष होऊ शकतो: $t_2$ ला $[s_1 \circ s_2]$ असे वाचल्यास $t_1 \sqsubseteq t_2$ मिळेल; पण $t_2$ ला $[r_1 \circ r_2]$ असे वाचल्यास $t_1 \sqsubseteq t_2$ असत्य मिळेल. हे घडू शकत नाही हे सिद्ध करण्याआधी, जिथे ते \emph{घडू शकते} असे एक उदाहरण पाहू.

\begin{defn}
\emph{कंसविरहित पदे} विगमनाधारित रीतीने अशी परिभाषित करा:
\begin{enumerate}
\item प्रत्येक अक्षर कंसविरहित पद आहे.
\item $s_1$ आणि $s_2$ कंसविरहित पदे असतील, तर $s_1 \circ s_2$ हे कंसविरहित पद आहे.
\item इतर काहीही कंसविरहित पद नाही.
\end{enumerate}
\end{defn}

कंसविरहित पदांची उदाहरणे म्हणजे $\mathrm{a}$, $\mathrm{b} \circ \mathrm{d}$ आणि $\mathrm{b} \circ \mathrm{a} \circ \mathrm{b}$. यांच्यासाठी “उपपद” ची व्याख्या वरच्या पद्धतीने केली, तर दुसरा खंड असा असेल:
\begin{center}
$t_2 = s_1 \circ s_2$ असेल, तर $t_1 \sqsubseteq t_2$ हे $t_1 = t_2$, $t_1 \sqsubseteq s_1$ किंवा $t_1 \sqsubseteq s_2$ असेल तेव्हाच आणि फक्त तेव्हाच सत्य असते.
\end{center}

आता $\mathrm{b} \circ \mathrm{a} \circ \mathrm{b}$ हे $s_1 \circ s_2$ या रूपाचे आहे, जिथे
\begin{align*}
s_1 & = \mathrm{b} \text{ आणि} & s_2 & = \mathrm{a} \circ \mathrm{b}.
\intertext{तेच पद $r_1 \circ r_2$ या रूपाचेही आहे, जिथे}
r_1 & = \mathrm{b} \circ \mathrm{a} \text{ आणि} & r_2 &= \mathrm{b}.
\end{align*}
आता $\mathrm{a} \circ \mathrm{b}$ हे $\mathrm{b} \circ \mathrm{a} \circ \mathrm{b}$ चे उपपद आहे का? पहिल्या वाचनानुसार उत्तर होय; दुसऱ्यानुसार नाही.

एखादे नीटस पद इतर नीटस पदांपासून कसे रचले आहे हे एकमेव रीतीने निश्चित होते, या गुणधर्माला \emph{एकमेव रीतीने वाचनीयता} म्हणतात. अशा विगमनाधारित वस्तूंवरील संबंधांच्या विगमनाधारित व्याख्या महत्त्वाच्या असल्यामुळे हा गुणधर्म सिद्ध केला पाहिजे.

\begin{prop}
$t$ हे नीटस पद आहे असे समजा. मग $t$ हे स्वतंत्र अक्षर आहे; किंवा अशी एकमेव रीतीने निश्चित होणारी नीटस पदे $s_1$, $s_2$ आहेत की~$t = [s_1 \circ s_2]$.
\end{prop}

\begin{proof}
$t$ स्वतंत्र अक्षर असेल, तर अट पूर्ण होते. म्हणून $t$ स्वतंत्र अक्षर नाही असे समजा. विगमनाधारित व्याख्येवरून मग $t$ हे काही नीटस पदे $s_1$ आणि~$s_2$ यांसाठी $[s_1 \circ s_2]$ या रूपाचे असले पाहिजे. ही पदे एकमेव रीतीने निश्चित होतात हे दाखवायचे आहे: $t = [r_1 \circ r_2]$ असेल, तर $s_1 = r_1$ आणि $s_2 = r_2$.

म्हणून $t = [s_1 \circ s_2]$ आणि $t = [r_1 \circ r_2]$ असेही समजा; येथे $s_1$, $s_2$, $r_1$, $r_2$ ही नीटस पदे आहेत. $s_1 = r_1$ आणि $s_2 = r_2$ दाखवायचे आहे. प्रथम $s_1$ आणि $r_1$ एकरूप असले पाहिजेत; तसे नसतील, तर त्यांपैकी एक दुसऱ्याचा अरिक्त उचित पूर्वखंड असेल. पण \olref[sti]{prop:initial} नुसार $s_1$ आणि~$r_1$ दोन्ही नीटस पदे असतील, तर हे अशक्य आहे. $s_1 = r_1$ असल्यावर $s_2 = r_2$ देखील स्पष्टपणे मिळते.
\end{proof}

फलनेही विगमनाधारित रीतीने परिभाषित करता येतात. उदाहरणार्थ, कोणत्याही नीटस पदातील आतल्या आत असलेल्या $[\dots]$ जोड्यांची कमाल खोली देणारे~$f$ फलन पुढीलप्रमाणे परिभाषित करता येते:

\begin{defn}
\ollabel{defn:depth} नीटस पदाची \emph{खोली}, $f(t)$, पुढीलप्रमाणे विगमनाधारित व्याख्येने दिली आहे:
\[
f(t) = \begin{cases}
0 & \text{ जर $t$ अक्षर असेल}\\
\max(f(s_1), f(s_2)) + 1 & \text{ जर $t = [s_1 \circ s_2]$ असेल.}
\end{cases}
\]
\end{defn}

उदाहरणार्थ,
\begin{align*}
f([\mathrm{a} \circ \mathrm{b}]) & =
\max(f(\mathrm{a}),f(\mathrm{b})) + 1 = \\
&= \max(0, 0) + 1 = 1, \text{ आणि}\\
f([[\mathrm{a} \circ \mathrm{b}] \circ \mathrm{c}]) & =
\max(f([\mathrm{a} \circ \mathrm{b}]), f(\mathrm{c})) + 1 = \\
& = \max(1,0) + 1 = 2.
\end{align*}

अर्थात, येथे $s_1$ आणि $s_2$ ही नीटस पदे आहेत असे मानतो. प्रत्येक नीटस पद स्वतंत्र अक्षर किंवा $[s_1 \circ s_2]$ या रूपाचे असते, हे तथ्य वापरतो. ते या रूपात एकाच रीतीने असू शकते, हेही पुन्हा महत्त्वाचे आहे. का ते पाहण्यासाठी आधी परिभाषित केलेल्या कंसविरहित पदांचा पुन्हा विचार करा. त्यांच्यासाठी यासारखी “व्याख्या” अशी असेल:
\[
g(t) =
\begin{cases}
0 & \text{ जर $t$ अक्षर असेल}\\
\max(g(s_1), g(s_2)) + 1 & \text{ जर $t = s_1 \circ s_2$ असेल.}
\end{cases}
\]
आता $\mathrm{a} \circ \mathrm{b} \circ \mathrm{c} \circ \mathrm{d}$ हे कंसविरहित पद पाहा. ते अनेक रीतींनी वाचता येते. उदाहरणार्थ, $s_1 \circ s_2$ या रूपात, जिथे
\begin{align*}
s_1 & = \mathrm{a} \text{ आणि} &
s_2 & = \mathrm{b} \circ \mathrm{c} \circ \mathrm{d},
\intertext{किंवा $r_1 \circ r_2$ या रूपात, जिथे}
r_1 & = \mathrm{a} \circ b \text{ आणि} &
r_2 &= \mathrm{c} \circ \mathrm{d}.
\end{align*}
पहिल्या वाचनानुसार $g$ मोजल्यास असे मिळेल:
\begin{align*}
g(s_1 \circ s_2) & =
\max(g(\mathrm{a}), g(\mathrm{b} \circ \mathrm{c} \circ \mathrm{d})) + 1 =\\ & =
\max(0,2) + 1 = 3
\intertext{तर दुसऱ्या वाचनानुसार असे मिळते}
g(r_1 \circ r_2) & =
\max(g(\mathrm{a} \circ \mathrm{b}), g(\mathrm{c} \circ \mathrm{d})) + 1=\\ &
= \max(1,1) + 1 = 2
\end{align*}
पण फलनाने नेहमी एकमेव मूल्य दिले पाहिजे. म्हणून~$g$ ची ही “व्याख्या” मुळात फलनच परिभाषित करत नाही.

\begin{prob}
$l$ फलनाची विगमनाधारित व्याख्या द्या; येथे $l(t)$ म्हणजे नीटस पद~$t$ मधील चिन्हांची संख्या.
\end{prob}

\begin{prob}
नीटस पद~$t$ वर संरचनात्मक विगमन वापरून $f(t) < l(t)$ सिद्ध करा. येथे $l(t)$ म्हणजे~$t$ मधील चिन्हांची संख्या आणि $f(t)$ म्हणजे \olref[mth][ind][rel]{defn:depth} मध्ये परिभाषित केलेली~$t$ ची खोली.
\end{prob}

\end{document}
